\chapter{Answering the research question}\label{ch:conclusion}
\section{The answer and its basis}
Answer the question stated in \LawChapterRef{ch:introduction}. Identify the
principal findings and the reasoning that connects them. Explain the original
contribution to legal understanding without expanding its scope beyond the
analysis completed.

\section{The boundaries of the conclusion}
State the limits that matter to the answer. Identify developments after the
legal cut-off that require attention, questions the evidence cannot settle and
the conditions under which the conclusion may apply elsewhere.

\section{Further research and practical use}
Describe the next research question that follows from the thesis and the work
needed to answer it. If the thesis has practical implications, identify the
people or institutions able to act on them and the information they would need.

\LawPrompt{Check that every claim of contribution in the abstract and
introduction is supported by the completed thesis, and that the conclusion
answers the question actually investigated.}
